\FigureLayoutDeclare{img/2011/liaoning_science/q20_fig1.png}{6.0cm}{0bp 0bp 82.99bp 23.997bp}

\chapter{2011年辽宁卷（理）}

\section{单选题}

\begin{problem}
若 \(a\) 为正实数，\(\i\) 为虚数单位，\(\left|\frac{a + \i}{\i}\right| = 2\)，则 \(a =\)（\quad）

\choices
    {\(2\)}
    {\(\sqrt{3}\)}
    {\(\sqrt{2}\)}
    {\(1\)}
\end{problem}

\begin{answer}
B
\end{answer}

\begin{solution}
因为
\[
\left|\frac{a+\i}{\i}\right|=|1-a\i|=\sqrt{1+a^{2}}=2
\]
所以 \(a^{2}=3\). 又因为 \(a>0\)，故 \(a=\sqrt{3}\)，选 B.
\end{solution}

\begin{problem}
已知 \(M\)，\(N\) 为集合 \(I\) 的非空真子集，且 \(M\)，\(N\) 不相等，若 \(N \cap \complement_I M = \varnothing\)，则 \(M \cup N =\)（\quad）

\choices
    {\(M\)}
    {\(N\)}
    {\(I\)}
    {\(\varnothing\)}
\end{problem}

\begin{answer}
A
\end{answer}

\begin{solution}
由 \(N\cap\complement_I M=\varnothing\)，可知 \(N\) 中不存在属于 \(I\) 且不属于 \(M\) 的元素，因此 \(N\subseteq M\). 题设又给出 \(M\ne N\)，所以 \(N\) 是 \(M\) 的真子集，从而 \(M\cup N=M\)，选 A.
\end{solution}

\begin{problem}
已知 \(F\) 是抛物线 \(y^{2} = x\) 的焦点，\(A\)，\(B\) 是该抛物线上的两点，\(|AF| + |BF| = 3\)，则线段 \(AB\) 的中点到 \(y\) 轴的距离为（\quad）

\choices
    {\( \frac{3}{4} \)}
    {\(1\)}
    {\( \frac{5}{4} \)}
    {\( \frac{7}{4} \)}
\end{problem}

\begin{answer}
C
\end{answer}

\begin{solution}
抛物线 \(y^{2}=x\) 可写为 \(y^{2}=4px\)，故 \(p=\frac14\)，焦点为 \(F\left(\frac14,0\right)\)，准线为 \(x=-\frac14\). 设 \(A\)，\(B\) 的横坐标分别为 \(x_A\)，\(x_B\). 由抛物线的定义，\(|AF|=x_A+\frac14\)，\(|BF|=x_B+\frac14\).
因此 \(x_A+x_B+\frac12=3\)，即 \(x_A+x_B=\frac52\). 线段 \(AB\) 的中点横坐标为 \(\frac{x_A+x_B}{2}=\frac54\)，它到 \(y\) 轴的距离为 \(\frac54\)，选 C.
\end{solution}

\begin{problem}
\(\triangle ABC\) 的三个内角 \(A\)，\(B\)，\(C\) 所对的边分别为 \(a\)，\(b\)，\(c\)，并且有 \(a \sin A \sin B + b \cos^2 A = \sqrt{2} a\)，则 \(\frac{b}{a} =\)（\quad）

\choices
    {\(2 \sqrt{3}\)}
    {\(2 \sqrt{2}\)}
    {\(\sqrt{3}\)}
    {\(\sqrt{2}\)}
\end{problem}

\begin{answer}
D
\end{answer}

\begin{solution}
由正弦定理，\(\frac ba=\frac{\sin B}{\sin A}\). 将其代入题设等式并除以 \(a\)，得
\[
\sin A\sin B+\frac{\sin B}{\sin A}\cos^{2}A=\sqrt2
\]
左边等于
\[
\sin B\left(\sin A+\frac{\cos^{2}A}{\sin A}\right)=\frac{\sin B}{\sin A}
\]
所以 \(\frac ba=\sqrt2\)，选 D.
\end{solution}

\begin{problem}
从 1， 2， 3， 4， 5 中任取 2 个不同的数，事件 \(A =\) “取到的 2 个数之和为偶数”，事件 \(B =\) “取到的 2 个数均为偶数”，则 \( P(B \mid A) = \)（\quad）

\choices
    {\(\frac{1}{8}\)}
    {\(\frac{1}{4}\)}
    {\(\frac{2}{5}\)}
    {\(\frac{1}{2}\)}
\end{problem}

\begin{answer}
B
\end{answer}

\begin{solution}
从 \(1\)，\(2\)，\(3\)，\(4\)，\(5\) 中任取两个数共有 \(\mathrm C_5^2=10\) 个等可能结果. 事件 \(A\) 发生时，两个数必须同奇偶：两个奇数的取法有 \(\mathrm C_3^2=3\) 个，两个偶数的取法有 \(\mathrm C_2^2=1\) 个，所以 \(P(A)=\frac4{10}\). 事件 \(B\) 只有取到 \(\{2,4\}\) 这一种结果，且 \(B\subset A\)，故
\[
P(B\mid A)=\frac{P(B)}{P(A)}=\frac{1/10}{4/10}=\frac14
\]
选 B.
\end{solution}

\begin{problem}
执行下面的程序框图，如果输入的 \(n\) 是 4，则输出的 \(p\) 是（\quad）

\bitmapfigure[max width=0.42\linewidth]{img/2011/liaoning_science/q6_fig1.png}

\choices
    {\(8\)}
    {\(5\)}
    {\(3\)}
    {\(2\)}
\end{problem}

\begin{answer}
C
\end{answer}

\begin{solution}
初始时 \(s=0\)，\(t=1\)，\(k=1\)，\(p=1\). 第一次循环后 \(s=1\)，\(t=1\)，\(k=2\)，\(p=2\)，此时 \(k<4\)，继续循环；第二次循环后 \(s=1\)，\(t=2\)，\(k=3\)，\(p=3\)，仍有 \(k<4\)，继续循环；第三次循环后 \(s=2\)，\(t=3\)，\(k=4\)，\(p=3\)，循环停止. 因此输出 \(p=3\)，选 C.
\end{solution}

\begin{problem}
设 \(\sin \left(\frac{\pi}{4} + \theta\right) = \frac{1}{3}\)，则 \(\sin 2\theta =\)（\quad）

\choices
    {\(-\frac{7}{9}\)}
    {\(-\frac{1}{9}\)}
    {\( \frac{1}{9} \)}
    {\( \frac{7}{9} \)}
\end{problem}

\begin{answer}
A
\end{answer}

\begin{solution}
令 \(\alpha=\frac\pi4+\theta\)，则 \(\sin\alpha=\frac13\). 因为
\[
\sin2\theta=\sin\left(2\alpha-\frac\pi2\right)=-\cos2\alpha=-(1-2\sin^{2}\alpha)
\]
所以 \(\sin2\theta=-1+\frac29=-\frac79\)，选 A.
\end{solution}

\begin{problem}
如图，四棱锥 \(S - ABCD\) 的底面为正方形，\(SD \perp\) 底面 \(ABCD\)，则下列结论中不正确的是（\quad）

\bitmapfigure[max width=0.34\linewidth]{img/2011/liaoning_science/q8_fig1.png}

\choices
    {\(AC \perp SB\)}
    {\(AB\parallel\) 平面 \(SCD\)}
    {\(SA\) 与平面 \(SBD\) 所成的角等于 \(SC\) 与平面 \(SBD\) 所成的角}
    {\(AB\) 与 \(SC\) 所成的角等于 \(DC\) 与 \(SA\) 所成的角}
\end{problem}

\begin{answer}
D
\end{answer}

\begin{solution}
设正方形边长为 \(1\)，取坐标 \(D=(0,0,0)\)，\(A=(1,0,0)\)，\(C=(0,1,0)\)，\(B=(1,1,0)\)，\(S=(0,0,h)\)，
其中 \(h>0\). 则 \(\overrightarrow{AC}=(-1,1,0)\)，\(\overrightarrow{SB}=(1,1,-h)\)，两者内积为 \(0\)，所以 \(AC\perp SB\). 又 \(AB\parallel DC\)，而 \(DC\) 在平面 \(SCD\) 内，所以 \(AB\parallel\) 平面 \(SCD\).

平面 \(SBD\) 的一个法向量为 \((1,-1,0)\). 向量 \(\overrightarrow{SA}=(1,0,-h)\) 与 \(\overrightarrow{SC}=(0,1,-h)\) 在该法向量上的投影绝对值都为 \(1\)，且两向量长度都为 \(\sqrt{1+h^{2}}\)，故 \(SA\) 与平面 \(SBD\) 的角等于 \(SC\) 与平面 \(SBD\) 的角.

另一方面，\(AB\) 与 \(SC\) 的夹角满足 \(\cos\angle(AB,SC)=\frac1{\sqrt{1+h^{2}}}>0\)，而 \(DC\perp SA\)，所以 \(DC\) 与 \(SA\) 的夹角为 \(90^\circ\)，两角不相等. 故不正确的是 D.
\end{solution}

\begin{problem}
设函数 \( f(x) = \begin{cases}
2^{1 - x}, & x \leqslant 1,\\
1 - \log_2 x, & x > 1,
\end{cases} \) 则满足 \( f(x) \leqslant 2 \) 的 \( x \) 取值范围是（\quad）

\choices
    {\([-1,2]\)}
    {\([0,2]\)}
    {\([1,+\infty)\)}
    {\([0,+\infty)\)}
\end{problem}

\begin{answer}
D
\end{answer}

\begin{solution}
当 \(x\leqslant1\) 时，\(2^{1-x}\leqslant2=2^1\)，因底数 \(2>1\)，所以 \(1-x\leqslant1\)，即 \(x\geqslant0\). 与 \(x\leqslant1\) 合并得 \(x\in[0,1]\).

当 \(x>1\) 时，\(1-\log_2x\leqslant2\) 等价于 \(\log_2x\geqslant-1\)，即 \(x\geqslant\frac12\)，与 \(x>1\) 合并仍为 \(x>1\). 因此解集为 \([0,+\infty)\)，选 D.
\end{solution}

\begin{problem}
若 \(\bs{a}\)，\(\bs{b}\)，\(\bs{c}\) 均为单位向量，且 \(\bs{a} \cdot \bs{b} = 0\)，\((\bs{a} - \bs{c}) \cdot (\bs{b} - \bs{c}) \leqslant 0\)，则 \(\lvert \bs{a} + \bs{b} - \bs{c} \rvert\) 的最大值为（\quad）

\choices
    {\(\sqrt{2} - 1\)}
    {\(1\)}
    {\(\sqrt{2}\)}
    {\(2\)}
\end{problem}

\begin{answer}
B
\end{answer}

\begin{solution}
由 \(|\bs{a}|=|\bs{b}|=|\bs{c}|=1\)，\(\bs{a}\cdot\bs{b}=0\)，以及
\[
(\bs{a}-\bs{c})\cdot(\bs{b}-\bs{c})=\bs{a}\cdot\bs{b}-\bs{a}\cdot\bs{c}-\bs{b}\cdot\bs{c}+\bs{c}\cdot\bs{c}
=1-(\bs{a}+\bs{b})\cdot\bs{c}\leqslant0
\]
得 \((\bs{a}+\bs{b})\cdot\bs{c}\geqslant1\). 令 \(\bs{u}=\bs{a}+\bs{b}\)，则
\(|\bs{u}|^2=|\bs{a}|^2+|\bs{b}|^2+2\bs{a}\cdot\bs{b}=2\)，从而
\[
|\bs{a}+\bs{b}-\bs{c}|^2=|\bs{u}-\bs{c}|^2=|\bs{u}|^2+|\bs{c}|^2-2\bs{u}\cdot\bs{c}=3-2\bs{u}\cdot\bs{c}\leqslant1
\]
故 \(|\bs{a}+\bs{b}-\bs{c}|\leqslant1\). 取 \(\bs{c}=\bs{a}\)，则 \(\bs{c}\) 仍为单位向量，且 \((\bs{a}-\bs{c})\cdot(\bs{b}-\bs{c})=0\leqslant0\)，\(|\bs{a}+\bs{b}-\bs{c}|=|\bs{b}|=1\). 所以最大值为 \(1\)，选 B.
\end{solution}

\begin{problem}
函数 \( f(x) \) 的定义域为 \( \R \)，\( f(-1) = 2 \)，对任意 \( x \in \R \)，\( f'(x) > 2 \)，则 \( f(x) > 2x + 4 \) 的解集为（\quad）

\choices
    {\((-1,1)\)}
    {\((-1,+\infty)\)}
    {\((-\infty,-1)\)}
    {\((-\infty, +\infty)\)}
\end{problem}

\begin{answer}
B
\end{answer}

\begin{solution}
令 \(g(x)=f(x)-2x-4\). 由 \(f(-1)=2\) 得 \(g(-1)=0\)，又
\[
g'(x)=f'(x)-2>0
\]
所以 \(g(x)\) 在 \(\R\) 上严格递增，从而 \(g(x)>0\) 当且仅当 \(x>-1\). 因此不等式的解集为 \((-1,+\infty)\)，选 B.
\end{solution}

\begin{problem}
已知球的直径 \(SC = 4\)，\(A\)，\(B\) 是该球球面上的两点，\(AB = \sqrt{3}\)，\(\angle ASC = \angle BSC = 30^\circ\)，则棱锥 \(S - ABC\) 的体积为（\quad）

\choices
    {\(3\sqrt{3}\)}
    {\(2\sqrt{3}\)}
    {\(\sqrt{3}\)}
    {\(1\)}
\end{problem}

\begin{answer}
C
\end{answer}

\begin{solution}
取 \(S\) 为原点，令 \(SC\) 为 \(x\) 轴正方向. 由半径为 \(2\) 的球面及 \(\angle ASC=\angle BSC=30^\circ\)，可取 \(C=(4,0,0)\)，\(A=(3,\sqrt3,0)\)，\(B=\left(3,\frac{\sqrt3}{2},\frac32\right)\).
这里 \(SA=SB=2\sqrt3\)，且
\[
AB^2=\left(\frac{\sqrt3}{2}\right)^2+\left(\frac32\right)^2=3
\]
所以
\[
V_{S-ABC}=\frac16\left|\det\begin{pmatrix}3&\sqrt3&0\\3&\frac{\sqrt3}{2}&\frac32\\4&0&0\end{pmatrix}\right|=\frac16\cdot6\sqrt3=\sqrt3
\]
选 C.
\end{solution}

\section{填空题}

\begin{problem}
已知点 \((2,3)\) 在双曲线 \(C: \frac{x^2}{a^2} - \frac{y^2}{b^2} = 1 \) （\(a > 0\)，\(b > 0\)）上，\(C\) 的焦距为 4，则它的离心率为\fillinblank{}.
\end{problem}

\begin{answer}
2
\end{answer}

\begin{solution}
由双曲线的焦距为 \(4\)，得 \(2c=4\)，所以 \(c=2\). 双曲线满足 \(c^{2}=a^{2}+b^{2}\)，故
\[
a^{2}+b^{2}=4
\]
又因为点 \((2,3)\) 在双曲线上，
\[
\frac4{a^{2}}-\frac9{b^{2}}=1
\]
令 \(A=a^{2}\)，则 \(b^{2}=4-A\)，代入上式得
\(\frac4A-\frac9{4-A}=1\)，\(\Longrightarrow\)，\(A^{2}-17A+16=0\).
所以 \(A=1\) 或 \(A=16\). 由于 \(a^{2}<c^{2}=4\)，只能取 \(a^{2}=1\)，从而 \(a=1\). 因此离心率为
\[
e=\frac ca=2
\]
\end{solution}

\begin{problem}
调查了某地若干户家庭的年收入 \(x\) （单位：万元） 和年饮食支出 \(y\) （单位：万元），调查显示年收入 \(x\) 与年饮食支出 \(y\) 具有线性相关关系. 并由调查数据得到 \(y\) 对 \(x\) 的回归直线方程： \(\hat{y} = 0.254x + 0.321\). 由回归直线方程可知，家庭收入每增加 1 万元，年饮食支出平均增加\fillinblank{} 万元.
\end{problem}

\begin{answer}
0.254
\end{answer}

\begin{solution}
回归直线 \(\hat y=0.254x+0.321\) 的斜率为 \(0.254\). 当家庭收入增加 \(1\) 万元时，回归值的变化量为
\[
\Delta\hat y=0.254\times1=0.254
\]
万元，所以年饮食支出平均增加 \(0.254\) 万元.
\end{solution}

\begin{problem}
一个正三棱柱的侧棱长和底面边长相等，体积为 \(2 \sqrt{3}\)，它的三视图中的俯视图如图所示，左视图是一个矩形，则这个矩形的面积是\fillinblank{}

\bitmapfigure[max width=0.34\linewidth]{img/2011/liaoning_science/q15_fig1.png}
\end{problem}

\begin{answer}
\(2\sqrt{3}\)
\end{answer}

\begin{solution}
设正三棱柱的底面边长和侧棱长都为 \(a\). 正三角形底面积为 \(\frac{\sqrt3}{4}a^{2}\)，所以由体积条件
\(\frac{\sqrt3}{4}a^{2}\cdot a=2\sqrt3\)，\(\Longrightarrow\)，\(a^{3}=8\).
从而 \(a=2\). 按图从左侧观察时，左视图矩形的水平边是正三角形在观察方向上的投影，即正三角形的高，长度为 \(\frac{\sqrt3}{2}a=\sqrt3\)；另一边为棱柱的侧棱长 \(a=2\). 因此矩形面积为
\[
2\cdot\sqrt3=2\sqrt3
\]
\end{solution}

\begin{problem}
已知函数 \( f(x) = A \tan (\omega x + \varphi) \)（\(\omega > 0\)，\(|\varphi| < \frac{\pi}{2}\)），\( y = f(x) \) 的部分图象如图，则 \( f\left(\frac{\pi}{24}\right) = \)\fillinblank{}.

\bitmapfigure[max width=0.34\linewidth]{img/2011/liaoning_science/q16_fig1.png}
\end{problem}

\begin{answer}
\(\sqrt{3}\)
\end{answer}

\begin{solution}
由图可知相邻两条竖直渐近线的横坐标之差为 \(\frac{\pi}{2}\). 正切函数相邻渐近线的自变量间隔为 \(\pi\)，故
\[
\frac{\pi}{\omega}=\frac{\pi}{2}\Longrightarrow \omega=2
\]
图象与 \(x\) 轴交于 \(x=\frac{3\pi}{8}\)，因此
\[
2\cdot\frac{3\pi}{8}+\varphi=k\pi
\]
结合 \(\lvert\varphi\rvert<\frac{\pi}{2}\)，可取 \(k=1\)，从而 \(\varphi=\frac{\pi}{4}\). 又图象经过 \((0,1)\)，所以 \(A\tan\frac{\pi}{4}=A=1\). 因此
\[
f\left(\frac{\pi}{24}\right)
=\tan\left(2\cdot\frac{\pi}{24}+\frac{\pi}{4}\right)
=\tan\frac{\pi}{3}
=\sqrt3
\]
\end{solution}

\section{解答题}

\begin{problem}
已知等差数列 \(\{a_{n}\}\) 满足 \(a_{2} = 0\)，\(a_{6} + a_{8} = -10\).
\begin{enumerate}
\item 求数列 \( \{a_{n}\} \) 的通项公式；
\item 求数列 \( \left\{\frac{a_{n}}{2^{n-1}}\right\} \) 的前 \(n\) 项和.
\end{enumerate}
\end{problem}

\begin{solution}
\begin{enumerate}
\item 设首项为 \(a_{1}\)，公差为 \(d\). 由 \(a_{2}=0\) 得
\[
a_{1}+d=0
\]
又
\[
a_{6}+a_{8}=(a_{1}+5d)+(a_{1}+7d)=2a_{1}+12d=-10
\]
由 \(a_{1}=-d\) 得 \(10d=-10\)，所以 \(d=-1\)，\(a_{1}=1\). 因此
\[
a_{n}=a_{1}+(n-1)d=1-(n-1)=2-n
\]
\item 令
\(T_{n}=\sum_{k=1}^{n}\frac{a_{k}}{2^{k-1}}\).
当 \(n=1\) 时，\(T_{1}=a_{1}=1=\frac1{2^{0}}\). 假设 \(T_{n-1}=\frac{n-1}{2^{n-2}}\)，则
\[
T_{n}=T_{n-1}+\frac{a_{n}}{2^{n-1}}
=\frac{n-1}{2^{n-2}}+\frac{2-n}{2^{n-1}}
=\frac{2n-2+2-n}{2^{n-1}}
=\frac{n}{2^{n-1}}
\]
由数学归纳法，对任意正整数 \(n\)，都有
\[
\sum_{k=1}^{n}\frac{a_{k}}{2^{k-1}}=\frac{n}{2^{n-1}}
\]
\end{enumerate}
\end{solution}

\begin{problem}
如图，四边形 \(ABCD\) 为正方形，\(PD \perp\) 平面 \(ABCD\)，\(PD \parallel QA\)，\(QA = AB = \frac{1}{2} PD\).
\begin{enumerate}
\item 证明：平面 \(PQC \perp\) 平面 \(DCQ\)；
\item 求二面角 \(Q - BP - C\) 的余弦值.
\end{enumerate}

\bitmapfigure[max width=0.38\linewidth]{img/2011/liaoning_science/q18_fig1.png}
\end{problem}

\begin{solution}
\begin{enumerate}
\item 建立空间直角坐标系，令
\(A=(0,0,0)\)，\(B=(1,0,0)\)，\(C=(1,1,0)\)，\(D=(0,1,0)\).
由 \(PD\perp\) 平面 \(ABCD\)、\(PD=2\)、\(QA=1\) 及 \(PD\parallel QA\)，可取
\(P=(0,1,2)\)，\(Q=(0,0,1)\).
于是 \(\overrightarrow{PQ}=(0,-1,-1)\)，\(\overrightarrow{PC}=(1,0,-2)\)，\(\overrightarrow{DC}=(1,0,0)\)，\(\overrightarrow{DQ}=(0,-1,1)\).
平面 \(PQC\) 的一个法向量为
\[
\overrightarrow{PQ}\times\overrightarrow{PC}=(2,-1,1)
\]
平面 \(DCQ\) 的一个法向量为
\[
\overrightarrow{DC}\times\overrightarrow{DQ}=(0,-1,-1)
\]
这两个法向量的内积为 \(0\)，所以平面 \(PQC\perp\) 平面 \(DCQ\).
\item 平面 \(QBP\) 与平面 \(CBP\) 的法向量可分别取
\[
\overrightarrow{BQ}\times\overrightarrow{BP}=(-1,1,-1)
\]
以及
\[
\overrightarrow{BC}\times\overrightarrow{BP}=(2,0,1)
\]
它们的内积为 \(-3\)，长度分别为 \(\sqrt3\) 和 \(\sqrt5\). 因而二面角 \(Q-BP-C\) 的余弦值为
\[
\frac{-3}{\sqrt3\sqrt5}=-\frac{\sqrt{15}}5
\]
\end{enumerate}
\end{solution}

\begin{problem}
某农场计划种植某种新作物，为此对这种作物的两个品种（分别称为品种甲和品种乙）进行田间试验，选取两大块地. 每大块地分成 \(n\) 小块地. 在总共 \(2n\) 小块地中，随机选 \(n\) 小块地种植品种甲，另外 \(n\) 小块地种植品种乙.
\begin{enumerate}
\item 假设 \(n = 4\)，在第一大块地中，种植品种甲的小块地的数目记为 \(X\)，求 \(X\) 的分布列和数学期望.
\item 试验时每大块地分成 \(8\) 小块，即 \(n = 8\)，试验结束后得到品种甲和品种乙在各小块地上的每公顷产量 （单位： \(\mathrm{kg/hm^2}\)）如下表：
\begin{center}
\begin{tabular}{c|cccccccc}
\hline
品种甲 & 403 & 397 & 390 & 404 & 388 & 400 & 412 & 406 \\
\hline
品种乙 & 419 & 403 & 412 & 418 & 408 & 423 & 400 & 413 \\
\hline
\end{tabular}
\end{center}
分别求品种甲和品种乙每公顷产量的样本平均数和样本方差；根据试验结果，你应该种植哪一品种？

附：样本数据 \(x_{1}\)，\(x_{2}\)，\(\cdots\)，\(x_{n}\) 的样本方差 \(s^{2} = \frac{1}{n}[(x_{1} - \overline{x})^{2} + (x_{2} - \overline{x})^{2} + \cdots + (x_{n} - \overline{x})^{2}]\)，其中 \(\overline{x}\) 为样本平均数.
\end{enumerate}
\end{problem}

\begin{solution}
\begin{enumerate}
\item 总共有 \(8\) 块地，其中有 \(4\) 块种植甲，随机抽取其中 \(4\) 块作为第一大地，所以 \(X\) 服从超几何分布. 对 \(k=0\)，\(1\)，\(2\)，\(3\)，\(4\)，
\[
P(X=k)=\frac{\mathrm C_{4}^{k}\mathrm C_{4}^{4-k}}{\mathrm C_{8}^{4}}
\]
由于 \(\mathrm C_{8}^{4}=70\)，所以 \(P(X=0)=\frac1{70}\)，\(P(X=1)=\frac8{35}\)，\(P(X=2)=\frac{18}{35}\). 又 \(P(X=3)=\frac8{35}\)，\(P(X=4)=\frac1{70}\).
其数学期望为
\[
E(X)=4\cdot\frac4{8}=2
\]
\item 品种甲的数据和为 \(3200\)，所以
\[
\overline{x}_{\text{甲}}=\frac{3200}{8}=400
\]
相对于平均数 \(400\) 的偏差为 \(3\)，\(-3\)，\(-10\)，\(4\)，\(-12\)，\(0\)，\(12\)，\(6\)，平方和为
\[
3^{2}+(-3)^{2}+(-10)^{2}+4^{2}+(-12)^{2}+0^{2}+12^{2}+6^{2}=458
\]
因此
\[
s_{\text{甲}}^{2}=\frac{458}{8}=57.25
\]
品种乙的数据和为 \(3296\)，所以
\[
\overline{x}_{\text{乙}}=\frac{3296}{8}=412
\]
相对于平均数 \(412\) 的偏差为 \(7\)，\(-9\)，\(0\)，\(6\)，\(-4\)，\(11\)，\(-12\)，\(1\)，平方和为
\[
7^{2}+(-9)^{2}+0^{2}+6^{2}+(-4)^{2}+11^{2}+(-12)^{2}+1^{2}=448
\]
因此
\[
s_{\text{乙}}^{2}=\frac{448}{8}=56
\]
品种乙的平均产量更高，且样本方差不大于品种甲，所以应种植品种乙.
\end{enumerate}
\end{solution}

\begin{problem}
已知函数 \( f(x) = \ln x - ax^2 + (2 - a)x \).
\begin{enumerate}
\item 讨论 \(f(x)\) 的单调性；
\item 设 \(a > 0\)，证明：当 \(0 < x < \frac{1}{a}\) 时，\(f\left(\frac{1}{a} + x\right) > f\left(\frac{1}{a} - x\right)\)；
\item 若函数 \( y = f(x) \) 的图象与 \( x \) 轴交于 \(A\)，\(B\) 两点，线段 \(AB\) 中点的横坐标为 \( x_0 \)，证明：\( f'(x_0) < 0 \).
\end{enumerate}
\end{problem}

\begin{solution}
\begin{enumerate}
\item 函数定义域为 \((0,+\infty)\)，且
\[
f'(x)=\frac1x-2ax+2-a
= -\frac{(2x+1)(ax-1)}{x}
\]
当 \(a\leqslant0\) 时，\(2x+1>0\)、\(ax-1<0\)，所以 \(f'(x)>0\)，函数在 \((0,+\infty)\) 上单调增加. 当 \(a>0\) 时，\(f'(x)>0\) 在 \(0<x<1/a\) 成立，\(f'(x)<0\) 在 \(x>1/a\) 成立，故函数在 \((0,1/a)\) 上单调增加，在 \((1/a,+\infty)\) 上单调减少.
\item 令 \(g(x)=f\left(\frac1a+x\right)-f\left(\frac1a-x\right)\)，并取 \(0\leqslant x<\frac1a\).
直接化简得
\[
g(x)=\ln\frac{1+ax}{1-ax}-2ax
\]
并且
\[
g'(x)=\frac{a}{1+ax}+\frac{a}{1-ax}-2a
=\frac{2a^{3}x^{2}}{1-a^{2}x^{2}}>0
\]
当 \(0<x<1/a\) 时成立. 又 \(g(0)=0\)，所以 \(g(x)>0\)，即
\[
f\left(\frac1a+x\right)>f\left(\frac1a-x\right)
\]
\item 设两个交点的横坐标为 \(x_{1}<x_{2}\)，则 \(f(x_{1})=f(x_{2})=0\). 若 \(a\leqslant0\)，由第（1）问函数严格递增，不可能有两个交点，所以 \(a>0\). 由第（1）问的单调性，必有
\[
0<x_{1}<\frac1a<x_{2}
\]
在第（2）问中取 \(x=\frac1a-x_{1}\)，得到
\[
f\left(\frac2a-x_{1}\right)>f(x_{1})=0
\]
其中 \(\frac2a-x_{1}>\frac1a\)，而 \(f\) 在 \((1/a,+\infty)\) 上严格递减，且 \(f(x_{2})=0\)，故
\[
\frac2a-x_{1}<x_{2}
\]
于是中点横坐标
\[
x_{0}=\frac{x_{1}+x_{2}}2>\frac1a
\]
因此 \(f'(x_{0})<0\).
\end{enumerate}
\end{solution}

\begin{problem}
如图，已知椭圆 \(C_1\) 的中心在原点 \(O\)，长轴左，右端点 \(M\)，\(N\) 在 \(x\) 轴上，椭圆 \(C_2\) 的短轴为 \(MN\)，且 \(C_1\)，\(C_2\) 的离心率都为 \(e\). 直线 \(l \perp MN\)，\(l\) 与 \(C_1\) 交于两点，与 \(C_2\) 交于两点，这四点按纵坐标从大到小依次为 \(A\)，\(B\)，\(C\)，\(D\).
\begin{enumerate}
\item 设 \(e = \frac{1}{2}\)，求 \(|BC|\) 与 \(|AD|\) 的比值.
\item 当 \(e\) 变化时，是否存在直线 \(l\)，使得 \(BO \parallel AN\)，并说明理由.
\end{enumerate}

\bitmapfigure[max width=0.38\linewidth]{img/2011/liaoning_science/q20_fig1.png}
\end{problem}

\begin{solution}
\begin{enumerate}
\item 设椭圆 \(C_{1}\) 的半长轴为 \(a\)，记
\[
r=\sqrt{1-e^{2}}
\]
因为 \(C_{1}\) 的半短轴为 \(ar\)，而 \(C_{2}\) 的短轴 \(MN\) 长为 \(2a\)，所以 \(C_{2}\) 的半长轴为 \(a/r\). 因而两椭圆方程可写为
\[
C_{1}:\frac{x^{2}}{a^{2}}+\frac{y^{2}}{a^{2}r^{2}}=1
\]
\[
C_{2}:\frac{x^{2}}{a^{2}}+\frac{r^{2}y^{2}}{a^{2}}=1
\]
设直线 \(l\) 的方程为 \(x=t\)，并令 \(u=\sqrt{a^{2}-t^{2}}\). 由 \(\lvert t\rvert<a\)，得到 \(A=\left(t,\frac{u}{r}\right)\)，\(B=(t,ru)\)，\(C=(t,-ru)\)，\(D=\left(t,-\frac{u}{r}\right)\).
所以
\[
\frac{\lvert BC\rvert}{\lvert AD\rvert}
=\frac{2ru}{2u/r}=r^{2}=1-e^{2}
\]
当 \(e=\frac12\) 时，该比值为 \(\frac34\).
\item 仍取上述坐标，设 \(N=(a,0)\). 由 \(BO\parallel AN\)，向量 \(\overrightarrow{BO}=(-t,-ru)\)，\(\overrightarrow{AN}=\left(a-t,-\frac{u}{r}\right)\) 平行，故其行列式为零：
\[
\det\begin{pmatrix}-t&-ru\\a-t&-u/r\end{pmatrix}
=u\left(\frac{t}{r}+r(a-t)\right)=0
\]
由于 \(u>0\)，得到
\(\frac{t}{r}+r(a-t)=0\)，\(\Longrightarrow\)，\(t=-\frac{ar^{2}}{1-r^{2}}=-\frac{a(1-e^{2})}{e^{2}}\).
要使直线 \(l\) 与两椭圆都各有两个交点，必须 \(\lvert t\rvert<a\)，即
\[
\frac{1-e^{2}}{e^{2}}<1
\quad\Longleftrightarrow\quad
e>\frac{\sqrt2}{2}
\]
结合椭圆离心率的范围 \(0<e<1\)，当 \(\frac{\sqrt2}{2}<e<1\) 时存在这样的直线，当 \(0<e\leqslant\frac{\sqrt2}{2}\) 时不存在.
\end{enumerate}
\end{solution}

\section{选考题：解答题}

\begin{problem}
如图，\(A\)，\(B\)，\(C\)，\(D\) 四点在同一圆上，\(AD\) 的延长线与 \(BC\) 的延长线交于 \(E\) 点，且 \(EC = ED\).
\begin{enumerate}
\item 证明： \(CD \parallel AB\)；
\item 延长 \(CD\) 到 \(F\)，延长 \(DC\) 到 \(G\)，使得 \(EF = EG\)，证明： \(A\)，\(B\)，\(G\)，\(F\) 四点共圆.
\end{enumerate}

\bitmapfigure[max width=0.34\linewidth]{img/2011/liaoning_science/q22_fig1.png}
\end{problem}

\begin{solution}
\begin{enumerate}
\item 由 \(EC=ED\)，三角形 \(ECD\) 为等腰三角形，所以
\[
\angle ECD=\angle CDE
\]
又 \(A\)，\(D\)，\(E\) 三点共线，且 \(A\)，\(B\)，\(C\)，\(D\) 四点共圆，于是（以下角均按有向角计算）
\[
\angle CDE=\angle CDA=\angle CBA=\angle EBA
\]
因此 \(\angle ECD=\angle EBA\). 由于 \(EC\) 与 \(EB\) 在同一直线上，得到
\[
CD\parallel AB
\]
\item 因为 \(ED=EC\)，点 \(E\) 在 \(DC\) 的垂直平分线上；又因为 \(EF=EG\)，点 \(E\) 在 \(FG\) 的垂直平分线上. 直线 \(F\)，\(D\)，\(C\)，\(G\) 共线，且两条垂直平分线都过 \(E\) 并垂直于该直线，所以它们是同一条直线，反射关于此直线把 \(D\)，\(C\) 互换，也把 \(F\)，\(G\) 互换.

由 \(AB\parallel DC\)，且圆 \(ABCD\) 的圆心在 \(DC\) 的垂直平分线上，上述反射还把 \(A\)，\(B\) 互换. 因而
\[
AF=BG
\]
又 \(AB\parallel FG\)，所以四边形 \(ABGF\) 为等腰梯形，其两腰所对的内角互补，故 \(A\)，\(B\)，\(G\)，\(F\) 四点共圆.
\end{enumerate}
\end{solution}

\begin{problem}
在平面直角坐标系 \(xOy\) 中，曲线 \(C_1\) 的参数方程为 \(\begin{cases}
x = \cos \varphi,\\
y = \sin \varphi.
\end{cases}\) （\(\varphi\) 为参数）. 曲线 \(C_2\) 的参数方程为 \(\begin{cases}
x = a\cos \varphi,\\
y = b\sin \varphi,
\end{cases}\) （\(a > b > 0\)，\(\varphi\) 为参数）. 在以 \(O\) 为极点，\(x\) 轴的正半轴为极轴的极坐标系中，射线 \(l: \theta = \alpha\) 与 \(C_1\)，\(C_2\) 各有一个交点. 当 \(\alpha = 0\) 时，这两个交点间的距离为 2，当 \(\alpha = \frac{\pi}{2}\) 时，这两个交点重合.
\begin{enumerate}
\item 分别说明 \(C_1\)，\(C_2\) 是什么曲线，并求出 \(a\) 与 \(b\) 的值；
\item 设当 \(\alpha = \frac{\pi}{4}\) 时，\(l\) 与 \(C_1\)，\(C_2\) 的交点分别为 \(A_1\)，\(B_1\)，当 \(\alpha = -\frac{\pi}{4}\) 时，\(l\) 与 \(C_1\)，\(C_2\) 的交点分别为 \(A_2\)，\(B_2\)，求四边形 \(A_1A_2B_2B_1\) 的面积.
\end{enumerate}
\end{problem}

\begin{solution}
\begin{enumerate}
\item 由参数方程 \(x=\cos\varphi\)，\(y=\sin\varphi\) 可得 \(x^{2}+y^{2}=1\)，所以 \(C_{1}\) 是以原点为圆心、半径为 \(1\) 的圆. 同理，
\[
\frac{x^{2}}{a^{2}}+\frac{y^{2}}{b^{2}}=1
\]
所以 \(C_{2}\) 是椭圆.

当 \(\alpha=0\) 时，射线为 \(x\) 轴正半轴，两个交点的横坐标分别为 \(1\) 和 \(a\)，由距离为 \(2\) 且 \(a>1\)，得 \(a-1=2\)，即 \(a=3\). 当 \(\alpha=\frac{\pi}{2}\) 时，射线为 \(y\) 轴正半轴，两个交点的纵坐标分别为 \(1\) 和 \(b\)，由两点重合得 \(b=1\).
\item 当 \(\alpha=\frac{\pi}{4}\) 时，令 \(A_{1}=(u,u)\)，由 \(A_{1}\in C_{1}\) 得
由 \(2u^{2}=1\) 及 \(u>0\)，得 \(u=\frac{\sqrt2}{2}\).
令 \(B_{1}=(v,v)\)，由 \(B_{1}\in C_{2}\) 及 \(a=3\)，\(b=1\) 得
\[
\frac{v^{2}}9+v^{2}=1\Longrightarrow v=\frac{3\sqrt{10}}{10}
\]
当 \(\alpha=-\frac{\pi}{4}\) 时，关于 \(x\) 轴对称，故 \(A_{2}=(u,-u)\)，\(B_{2}=(v,-v)\).
四边形 \(A_{1}A_{2}B_{2}B_{1}\) 的两条平行边长度分别为 \(2u\)，\(2v\)，两平行边间的距离为 \(v-u\). 所以面积为
\[
\frac{2u+2v}{2}(v-u)=v^{2}-u^{2}
=\frac9{10}-\frac12=\frac25
\]
\end{enumerate}
\end{solution}

\begin{problem}
已知函数 \(f(x) = |x - 2| - |x - 5|\).
\begin{enumerate}
\item 证明： \(-3 \leqslant f(x) \leqslant 3\).
\item 求不等式 \( f(x) \geqslant x^{2} - 8x + 15 \) 的解集.
\end{enumerate}
\end{problem}

\begin{solution}
\begin{enumerate}
\item 由绝对值三角不等式，
\[
\bigl||x-2|-|x-5|\bigr|
\leqslant |(x-2)-(x-5)|=3
\]
所以 \(-3\leqslant f(x)\leqslant3\).
\item 按 \(x\) 与 \(2,5\) 的大小关系分段：
\[
f(x)=
\begin{cases}
-3,&x\leqslant2,\\
2x-7,&2<x<5,\\
3,&x\geqslant5.
\end{cases}
\]
当 \(x\leqslant2\) 时，
\[
-3\geqslant x^2-8x+15
\quad\Longleftrightarrow\quad
x^2-8x+18\leqslant0
\]
该二次式判别式为 \(64-72=-8<0\)，恒为正，无解.

当 \(2<x<5\) 时，
\[
2x-7\geqslant x^2-8x+15
\quad\Longleftrightarrow\quad
x^2-10x+22\leqslant0
\quad\Longleftrightarrow\quad
(x-5)^2\leqslant3
\]
与 \(2<x<5\) 合并得 \(5-\sqrt3\leqslant x<5\).

当 \(x\geqslant5\) 时，
\[
3\geqslant x^2-8x+15
\quad\Longleftrightarrow\quad
x^2-8x+12\leqslant0
\quad\Longleftrightarrow\quad
2\leqslant x\leqslant6
\]
与 \(x\geqslant5\) 合并得 \(5\leqslant x\leqslant6\). 因此解集为
\[
[5-\sqrt3,6]
\]
\end{enumerate}
\end{solution}
